\section{The free directions and the arity-two equation}

\begin{definition}[Finite counterterms along $\Delta$ and $P_b$]
\label{def:c11-free}
Write the solutions of the arity-one equations along $\Delta$ and along $P_b$ as
\[
\begin{aligned}
\hbar F_1^{(1)}(\Delta)=\ &4\,\text{units}\cdot(\text{field-counting functional})+\Big[-\tfrac23\int(A,A^*)+\tfrac13\int(\lambda^*,\lambda)+\tfrac43\int(c^*,c)\Big]\,\text{units}\\
&+\theta\,\text{units}\int\big(2E(A|_{W_{\mathbb R}^\perp})+2A|_{W_{\mathbb R}^\perp},A^*|_{W_{\mathbb R}^\perp}\big)-\theta_\lambda\,\text{units}\int\big(2E(\lambda)+3\lambda,\lambda^*\big)\\
&+\theta_V\,\text{units}\Big[\int\big(2E(A|_V)+2A|_V,A^*|_V\big)-\int\big(2E(c),c^*\big)\Big]+\text{units}\cdot Q\Big(\kappa_V\int(\iota_xA|_V,c^*)+\kappa_\lambda\int(\rho(x)\lambda^*,\lambda^*)\Big)+(\text{cubic and higher}),\\
\hbar F_1^{(1)}(P_b)=\ &\theta_P\,\text{units}\int\big(\partial_b(A|_{W_{\mathbb R}^\perp}),A^*|_{W_{\mathbb R}^\perp}\big)-\theta_{P,\lambda}\,\text{units}\int(\partial_b\lambda,\lambda^*)+\theta_{P,V}\,\text{units}\Big[\int\big(\partial_b(A|_V),A^*|_V\big)-\int(\partial_bc,c^*)\Big]\\
&+Q\big(\text{a combination of }\textstyle\int(\iota_bA|_V,c^*)\text{ and }\int(\rho(b)\lambda^*,\lambda^*)\big)+(\text{cubic and higher}),
\end{aligned}
\]
with numbers $\theta,\theta_\lambda,\theta_V,\kappa_V,\kappa_\lambda,\theta_P,\theta_{P,\lambda},\theta_{P,V}$; $\theta$ is the coefficient of \eqref{eq:c9-Fperp} and $\theta_P$ that of Theorem \ref{thm:c9-P1}. The three terms with $\theta$, $\theta_\lambda$, $\theta_V$ are the sector parts of $S^{(1)}(\Delta)$ in the sectors of $A|_{W_{\mathbb R}^\perp}$, of $\lambda$ and of $(A|_V,c)$ times these numbers, and those with $\theta_P$, $\theta_{P,\lambda}$, $\theta_{P,V}$ the sector parts of $S^{(1)}(P_b)$ (Definition \ref{def:c11-kernels}). The scheme is translation covariant if the position part along $b$ of $\hbar F_1^{(1)}(\Delta)$ is $2\hbar F_1^{(1)}(P_b)$; then $\theta_P=\theta$, $\theta_{P,\lambda}=\theta_\lambda$, $\theta_{P,V}=\theta_V$.
\end{definition}

\begin{proposition}[Completeness of the quadratic directions]
\label{prop:c11-tc0}
Let $\mathcal F_{-1}$ and $\mathcal F_{-2}$ be the quadratic local functionals of degree $-1$ and $-2$, of Weyl weight $4$ with the position at weight $-1$, at most linear in the position and invariant under $\mathrm{Spin}(4)\times\mathrm{Spin}(6)$. On the subspaces invariant under the reflection of Theorem \ref{thm:c9-quadratic} (iii),
\[
\frac{\ker\big(Q|_{\mathcal F_{-1}}\big)}{Q\,\mathcal F_{-2}}=\operatorname{span}\big\{\text{the three sector parts of }S^{(1)}(\Delta)\big\},\qquad Q\,\mathcal F_{-2}=\operatorname{span}\Big\{Q\int(\iota_xA|_V,c^*),\ Q\int(\rho(x)\lambda^*,\lambda^*)\Big\}.
\]
Without the reflection the quotient has one further class, $\int(\lambda^*,\rho(e_1e_2e_3e_4)\lambda)$, with $e_1,\dots,e_4$ an oriented orthonormal basis of $V$.
\end{proposition}

\begin{proof}
By Schur's lemma, $\mathcal F_{-1}$ is spanned by $\int(A^*|_V,D\,A|_V)$ with the structures $D$ built from $1$, $x\cdot\partial$, $x_\mu\partial_\nu$, $x_\nu\partial_\mu$ on the vector index and $\epsilon_{\mu\nu\kappa\tau}x^\kappa\partial^\tau$; by $\int(A^*|_{W_{\mathbb R}^\perp},D\,A|_{W_{\mathbb R}^\perp})$ with $D\in\{1,x\cdot\partial\}$; by $\int(\lambda^*,D\lambda)$ with $D\in\{1,\ x\cdot\partial,\ \rho(x)\rho(\partial)\}$ times $\{1,\rho(e_1e_2e_3e_4)\}$; and by $\int(c^*,Dc)$ with $D\in\{1,x\cdot\partial\}$. $\mathcal F_{-2}$ is spanned by $\int x^\mu(A_\mu,c^*)$, $\int(\lambda^*,\rho(x)\lambda^*)$ and $\int(\lambda^*,\rho(e_1e_2e_3e_4)\rho(x)\lambda^*)$. The structures with $\epsilon$ or with one factor $\rho(e_1e_2e_3e_4)$ change sign under the reflection. With $Q$ of Proposition \ref{prop:c1-Q}, the ranks on the plane-wave normal forms, computed exactly, are, in the sectors of $A|_{W_{\mathbb R}^\perp}$, of $\lambda$ and of $(A|_V,c)$: $\dim\mathcal F_{-1}=2,3,6$, $\dim\ker Q=1,2,2$, $\dim Q\mathcal F_{-2}=0,1,1$ for the reflection-invariant subspaces. The kernel is spanned by the sector part of $A|_{W_{\mathbb R}^\perp}$; by the sector part of $\lambda$ and $Q\int(\rho(x)\lambda^*,\lambda^*)$; by the sector part of $(A|_V,c)$ and $Q\int(\iota_xA|_V,c^*)$. Without the reflection, $\dim\mathcal F_{-1}=15$, the $\epsilon$-structure is not $Q$-closed, $\int(\lambda^*,\rho(e_1e_2e_3e_4)\lambda)$ is closed and not exact, and $\int(\lambda^*,\rho(e_1e_2e_3e_4)\rho(x)\lambda^*)=0$.
\end{proof}

\begin{proposition}[What the arity-one equations fix]
\label{prop:c11-arity1}
\begin{itemize}
\item (i) Along $\Delta$ the quadratic part of \eqref{eq:c9-arity1} fixes the translation-invariant part $\big[-\tfrac23\int(A,A^*)+\tfrac13\int(\lambda^*,\lambda)+\tfrac43\int(c^*,c)\big]$ units (Theorem \ref{thm:c9-quadratic} (iii)) and leaves $\theta$, $\theta_\lambda$, $\theta_V$, $\kappa_V$, $\kappa_\lambda$ free. Along $P_b$ the quadratic part of the arity-one obstruction vanishes (Proposition \ref{prop:c11-position} (i)), so the quadratic part of $\hbar F_1^{(1)}(P_b)$ is $Q$-closed and every coefficient of Definition \ref{def:c11-free} along $P_b$ is free; the translation-covariant choice, $\hbar F_1^{(1)}(P_b)$ equal to half the position part along $b$ of $\hbar F_1^{(1)}(\Delta)$, solves the equation along $P_b$ at every order in the fields.
\item (ii) Two solutions along $\Delta$ differ by an $s$-closed local functional of degree $-1$ and Weyl weight $4$. The common shift of $\theta$, $\theta_\lambda$, $\theta_V$ by $t$ adds $t$ units times $S^{(1)}(\Delta)$, which is $s$-closed (Lemma \ref{lem:c9-SDelta} (iii)); the directions $Q\int(\iota_xA|_V,c^*)$ and $Q\int(\rho(x)\lambda^*,\lambda^*)$ complete to $s\int(\iota_xA|_V,c^*)$ and $s\int(\rho(x)\lambda^*,\lambda^*)$, and an $s$-exact direction $sG$ contributes $s\{G,O\}$ to $D_1(O)$ on cocycles, zero on cohomology.
\item (iii) Since $\{I,S^{(1)}(\Delta)\}=0$ while the brackets of $I$ with the sector parts are nonzero, the cubic part of \eqref{eq:c9-arity1} depends on the sector coefficients through $(\theta-\theta_V)$ times the bracket of the cubic part of $I$ with the sector part of $A|_{W_{\mathbb R}^\perp}$ plus $(\theta_\lambda-\theta_V)$ times its bracket with the sector part of $\lambda$. It is solvable exactly when this sum is congruent to the cubic part of the composite part minus the bracket of the cubic part of $I$ with the translation-invariant part, modulo the images under $Q$ of the cubic local functionals of degree $-1$. If the classes of the two brackets modulo the images under $Q$ of the cubic local functionals of degree $-1$, Weyl weight $4$ and degree at most one in the position are linearly independent, the cubic equation fixes $\theta-\theta_V$ and $\theta_\lambda-\theta_V$.
\end{itemize}
\end{proposition}

\begin{proof}
(i) Theorem \ref{thm:c9-quadratic} (iii) along $\Delta$, Proposition \ref{prop:c11-position} (i) along $P_b$. For the covariant choice: $s$ maps letters to letters and carries no position, so it commutes with taking the position part, and the position part of the arity-one obstruction along $\Delta$ is twice the obstruction along $P_b$ by Proposition \ref{prop:c11-position} (i), the term $4$ units times $s$ of the field-counting functional having no position part. (ii) Theorem \ref{thm:c8-obstruction} (iii) restricted to arity one along $\Delta$, and Lemma \ref{lem:c10-deriv}: $\{sG,O\}=-s\{G,O\}$ for $sO=0$ and $G$ of degree $-2$. (iii) $s$ of the sector terms is the sum of $\theta$, $\theta_\lambda$, $\theta_V$ units times the brackets of $I$ with the sector parts, whose sum over the three sectors vanishes; at cubic order $s$ of the solution is $Q$ of its cubic part plus the bracket of the cubic part of $I$ with its quadratic part.
\end{proof}

\begin{theorem}[The position part of $D_1$ is exact]
\label{thm:c11-arity2}
Assume a one-loop quantisation through arity two along $\Delta$ and $P_b$, that is, \eqref{eq:c8-qme} modulo $\hbar^2$ in its components of arity at most two along these generators, and local representatives through arity two along $(\Delta,P_b)$ on a space $\mathfrak F$ of local functionals stable under $\Lambda$, with $P_1$ translation covariant, $P_1(\int\omega\,D)=\int\omega\,P_1(D)$ for a map $P_1$ of densities commuting with the derivatives; this holds when $\hbar F_1^{(1)}(P_b)$ is translation invariant. Then for every $O=\int\omega\,D\in\mathfrak F$ with $sO=0$
\[
\big(\text{position part of }D_1\text{ along }b\big)(O)=\big[D_1,\{S^{(1)}(P_b),-\}\big](O)=\pm\,s\Big(\hbar L_1(O)^{(2)}(\Delta,P_b)\Big):
\]
the position part of $D_1$ is $s$-exact on cocycles, with the arity-two component of $O$ as primitive.
\end{theorem}

\begin{proof}
By Corollary \ref{cor:c8-homotopy} (ii), $[\Lambda^{(1)}(\Delta),\Lambda^{(1)}(P_b)]+\Lambda^{(1)}([\Delta,P_b])$ is a graded commutator of $\Lambda^{(0)}$ with an operator built from $\Lambda^{(2)}(\Delta,P_b)$. With $[\Delta,P_b]=-2P_b$ (Remark \ref{rem:c2-dilatation}) and $S^{(2)}(\Delta,P_b)=0$ (Definition \ref{def:c2-S2}), the order-$\hbar^0$ part of that operator vanishes and the order-$\hbar$ part of the identity reads
\[
\big[D_1,\{S^{(1)}(P_b),-\}\big]+\big[v.,P_1\big]-2P_1=\big[s,H_1\big],
\]
with $H_1$ built from $\hbar L_1^{(2)}(\Delta,P_b)$. For translation-covariant $P_1$, $v.(P_1O)=\int\omega\,2(x\cdot\partial+w_D+1)P_1(D)$ and $P_1(v.O)=\int\omega\,\big(2x^\mu P_1(\partial_\mu D)+2w_DP_1(D)\big)$, so $[v.,P_1]=2P_1$. Lemma \ref{lem:c11-canonical} identifies $[D_1,\{S^{(1)}(P_b),-\}]$ with the position part, and on a cocycle $[s,H_1]O=s(H_1O)$.
\end{proof}

\begin{lemma}[Position part of a bracket]
\label{lem:c11-bracketpos}
For a local functional $G=\int(x^\mu g_\mu+g_0)$, the position part along $b$ of $\{G,-\}$ is $\{\int b^\mu g_\mu,-\}$.
\end{lemma}

\begin{proof}
The bracket contracts a slot of $g_\mu$ at $y$ with a letter $\partial^\alpha\varphi$ of $D$ at $x$ through $K_0$, which gives $\partial^\alpha_x(x^\mu\cdots)$; the term with no derivative on $x^\mu$ is $x^\mu$ times the bracket with $\int g_\mu$, and the others carry no position.
\end{proof}

\begin{corollary}[The composite term along $P_b$ and the translation component]
\label{cor:c11-theta}
Under the hypotheses of Theorem \ref{thm:c11-arity2} and of Definition \ref{def:c10-hyp}, on every cocycle $O\in\mathfrak F$:
\begin{itemize}
\item (i) twice the composite term of $O$ along $P_b$ minus the bracket of $O$ with the position part along $b$ of $\hbar F_1^{(1)}(\Delta)$ is $s$-exact; this position part is
\[
2\theta\,\text{units}\int\big(\partial_b(A|_{W_{\mathbb R}^\perp}),A^*|_{W_{\mathbb R}^\perp}\big)-2\theta_\lambda\,\text{units}\int(\partial_b\lambda,\lambda^*)+2\theta_V\,\text{units}\Big[\int\big(\partial_b(A|_V),A^*|_V\big)-\int(\partial_bc,c^*)\Big]+\text{units}\cdot Q\Big(\kappa_V\int(\iota_bA|_V,c^*)+\kappa_\lambda\int(\rho(b)\lambda^*,\lambda^*)\Big)
\]
plus the position parts of the higher terms;
\item (ii) if $\theta=\theta_\lambda=\theta_V$, the composite term along $P_b$ induces $\theta$ units times $\{S^{(1)}(P_b),-\}$ on the cohomology of $\mathfrak F$, and on every family in which $\{S^{(1)}(P_b),O\}$ is not exact for some cocycle the composite term along $P_b$ determines $\theta$;
\item (iii) in the translation-covariant scheme the position part of $D_1$ is $2P_1$, and the arity-two equation along $(\Delta,P_b)$ is equivalent to $P_1\cong0$ on cohomology: $P_1(O)=\pm\tfrac12\,s\big(\hbar L_1(O)^{(2)}(\Delta,P_b)\big)$;
\item (iv) if $\theta_P=\theta_{P,\lambda}=\theta_{P,V}$, then $P_1$ induces $(\theta-\theta_P)$ units times $\{S^{(1)}(P_b),-\}$ on cohomology, and the finite part $\hbar\psi'=\tfrac a2\{S^{(1)}(\Delta),-\}$, $a=(\theta-\theta_P)$ units, of Proposition \ref{prop:c9-scheme} gives $P_1\cong0$ on every family at once, leaves $D_1$, $M_1$, $R_1$ unchanged and changes the components along $K_b$, $\Sigma$ and $\Sigma^*$ by $a$, $-\tfrac a2$ and $\tfrac a2$ times the classical transformation. If the differences of the coefficients along $P_b$ and along $\Delta$ depend on the sector, a multiple of $\{S^{(1)}(\Delta),-\}$ removes $P_1$ on the families with letters of one sector, and the translation-covariant scheme removes it on every family.
\end{itemize}
\end{corollary}

\begin{proof}
(i) By Corollary \ref{cor:c11-D1local} and Proposition \ref{prop:c11-position} (i), the position part of $D_1$ is twice the sum over the pairs of the composite pair kernels along $P_b$ minus the bracket with the position part of $\hbar F_1^{(1)}(\Delta)$, by Lemma \ref{lem:c11-bracketpos}; the coefficients of $\log\varepsilon$ and the bracket with the field-counting functional commute with the translations. Theorem \ref{thm:c11-arity2} gives exactness, and the position part of each sector part of $S^{(1)}(\Delta)$ is twice the sector part of $S^{(1)}(P_b)$ (Definition \ref{def:c11-kernels}). (ii) For equal coefficients the $\theta$-terms are $2\theta$ units times $S^{(1)}(P_b)$, which is $s$-closed, and the other directions are $s$-exact after completion, with $\{sG,O\}=-s\{G,O\}$. (iii) In the covariant scheme the position part of $\hbar F_1^{(1)}(\Delta)$ is $2\hbar F_1^{(1)}(P_b)$, and $P_1$ is the sum of the composite pair kernels along $P_b$ minus the bracket with $\hbar F_1^{(1)}(P_b)$ (Corollary \ref{cor:c11-D1local}). (iv) By (ii), $P_1\cong\theta$ units times $\{S^{(1)}(P_b),-\}$ minus $\theta_P$ units times the same; Proposition \ref{prop:c9-scheme} with $\lambda_{P_b}=-2$. For sector-dependent differences $P_1$ induces the sum over the sectors of the differences times the brackets with the sector parts of $S^{(1)}(P_b)$, plus the completions, which is a multiple of $\{S^{(1)}(P_b),-\}$ on the cocycles with letters of one sector.
\end{proof}

\section{The dilatation identities in kernel form}

Corollary \ref{cor:c10-P7} assumes that on cohomology the composite term minus the bracket with the finite counterterm beyond the field-counting functional induces $-\tfrac13$ units times $\{S^{(1)}(\Delta),-\}$. By Corollary \ref{cor:c11-D1local} this is a statement about the composite pair kernels, which splits into a statement along $P_b$ and a statement about the translation-invariant part.

\begin{lemma}[One-body terms]
\label{lem:c11-onebody}
For $\theta=\theta_\lambda=\theta_V$, $\kappa_V=\kappa_\lambda=0$ and the quadratic part of $\hbar F_1^{(1)}(\Delta)$ of Definition \ref{def:c11-free}, the translation-invariant part of $-\big\{\hbar F_1^{(1)}(\Delta)-4\,\text{units}\cdot(\text{field-counting functional}),X\big\}$ on a field letter $X=\partial^\alpha\varphi$ is $(c_\varphi-2\theta\,w_X)$ units times $X$, with $c_\varphi=\tfrac23$ for $A|_V$ and $A|_{W_{\mathbb R}^\perp}$, $\tfrac13$ for $\lambda$ and $\tfrac43$ for $c$, and its position part along $b$ is $-2\theta$ units times $\partial_bX$. At $\theta=-\tfrac13$ the translation-invariant part is $\tfrac23(2+|\alpha|)$ units times $X$ for every field slot and for $c$.
\end{lemma}

\begin{proof}
$\{\int(A,A^*),X\}=X$, $\{\int(\lambda^*,\lambda),\lambda\}=-\lambda$ and $\{\int(c^*,c),c\}=-c$ (Lemma \ref{lem:c4-count}), and the bracket of the sector part with a letter of its sector is $v.X$, with translation-invariant part $2w_XX$ and position part $2\partial_bX$ (Lemma \ref{lem:c11-canonical}). At $\theta=-\tfrac13$: $\tfrac23+\tfrac23(1+|\alpha|)$, $\tfrac13+\tfrac23(\tfrac32+|\alpha|)$ and $\tfrac43+\tfrac23|\alpha|$ are all $\tfrac23(2+|\alpha|)$.
\end{proof}

\begin{definition}[The identities in kernel form]
\label{def:c11-P7}
Let $\mathfrak F$ be a family of gauge-invariant cocycles of field letters, stable under $\Lambda$, satisfying (ii) and (iii) of Definition \ref{def:c10-hyp}. The letter count is the derivation that multiplies a monomial $\int\omega\,X_1\cdots X_n$ of letters $X_a=\partial^{\alpha_a}\varphi_a$, $\varphi_a\in\{A|_V,A|_{W_{\mathbb R}^\perp},\lambda\}$, by $\sum_a(2+|\alpha_a|)$; it is defined on canonical forms, and its commutator with $\{S^{(1)}(P_b),-\}$ is $\{S^{(1)}(P_b),-\}$. The three identities on $\mathfrak F$ are:
\begin{itemize}
\item (a) the cubic part of \eqref{eq:c9-arity1} forces $\theta=\theta_\lambda=\theta_V$;
\item (b) the sum over the pairs, with the Koszul signs of \eqref{eq:c10-decomp}, of the composite pair kernels along $P_b$ induces $-\tfrac13$ units times $\{S^{(1)}(P_b),-\}$ on the cohomology of $\mathfrak F$;
\item (c) the sum over the pairs of the translation-invariant parts of the composite pair kernels along $\Delta$ induces $-\tfrac23$ units times the letter count on the cohomology of $\mathfrak F$.
\end{itemize}
\end{definition}

\begin{remark}[Dependences among the identities]
\label{rem:c11-P7dep}
(1) By Proposition \ref{prop:c11-position} (ii), pair by pair the commutator of the translation-invariant part of the composite pair kernel along $\Delta$ with $\{S^{(1)}(P_b),-\}$ is twice the composite pair kernel along $P_b$, and the commutator of the letter count with $\{S^{(1)}(P_b),-\}$ is $\{S^{(1)}(P_b),-\}$; hence (c) on $\mathfrak F$ and on $\{S^{(1)}(P_b),\mathfrak F\}$ gives (b) along $b$, and the content of (b) beyond (c) lies in the directions $b$ with $\{S^{(1)}(P_b),\mathfrak F\}$ outside $\mathfrak F$. (2) Identity (a) follows from (b) on three families with letters of one sector each, under the hypotheses of Theorem \ref{thm:c11-arity2} (Proposition \ref{prop:c11-theta-sector}); the cubic arity-one equation is then a consistency condition. (3) Identities (b) and (c) are identities between graph sums at the regulator of Definition \ref{def:c8-PT}; the numbers $-\tfrac13$ and $-\tfrac23$ are finite parts of special-edge cycles, which depend on the cutoff shape (Remark \ref{rem:c10-scheme1} (d), Proposition \ref{prop:c10-shape}). The statement independent of the scheme is the spectrum of $D_1$ at the coefficients fixed by the arity-two equation (Remark \ref{rem:c10-scheme1}).
\end{remark}

\begin{theorem}[Consequences of the identities]
\label{thm:c11-P7}
Assume (a), (b) and (c) of Definition \ref{def:c11-P7} on $\mathfrak F$.
\begin{itemize}
\item (i) $D_1$ minus four times the map of Theorem \ref{thm:c10-brst} induces $-(\tfrac13+\theta)$ units times $\{S^{(1)}(\Delta),-\}$ on the cohomology of $\mathfrak F$: the assumption of Corollary \ref{cor:c10-P7}, for the solution with sector coefficient $\theta$.
\item (ii) Under the hypotheses of Theorem \ref{thm:c11-arity2}, and if $\{S^{(1)}(P_b),O\}$ is not exact for some cocycle $O\in\mathfrak F$, $\theta=-\tfrac13$, and $D_1$ induces four times the map of Theorem \ref{thm:c10-brst}; its spectrum on $\mathfrak F$ in a fixed classical Weyl weight is twice that of the one-loop dilatation coefficients, $-8$ units times the eigenvalues of $\sum_lH_{l,l+1}$ of [B03c, (1.1)--(1.3)] on planar families (Corollary \ref{cor:c10-beisert}).
\item (iii) In the translation-covariant scheme $P_1\cong0$ on $\mathfrak F$; otherwise Corollary \ref{cor:c11-theta} (iv) applies.
\end{itemize}
Conversely, if $D_1$ minus four times the map induces $-(\tfrac13+\theta)$ units times $\{S^{(1)}(\Delta),-\}$ for equal sector coefficients, then (b) and (c) hold on $\mathfrak F$.
\end{theorem}

\begin{proof}
(i) By Corollary \ref{cor:c11-D1local}, $D_1$ minus four times the map is the sum of the composite pair kernels along $\Delta$ minus the bracket with $\hbar F_1^{(1)}(\Delta)-4$ units times the field-counting functional. Its position part along $b$ is, by (b) and Lemma \ref{lem:c11-onebody}, $2\cdot(-\tfrac13)\partial_b-2\theta\,\partial_b=-2(\tfrac13+\theta)\partial_b$ units. Its translation-invariant part is, by (c) and Lemma \ref{lem:c11-onebody}, $-\tfrac23$ units times the letter count plus $\sum_a(c_{\varphi_a}-2\theta w_{X_a})$ units, which is $-2(\tfrac13+\theta)w_O$ units, by $\sum_a(c_{\varphi_a}+\tfrac23w_{X_a})=\tfrac23\sum_a(2+|\alpha_a|)$. Together they give $-(\tfrac13+\theta)(2x\cdot\partial+2w_O)$ units, which is $-(\tfrac13+\theta)$ units times $v.$ (Lemma \ref{lem:c11-canonical}). (ii) By Theorem \ref{thm:c11-arity2} the position part of $D_1$, which is $-2(\tfrac13+\theta)$ units times $\partial_b$ on cohomology, is exact, and $\partial_bO$ is not exact for that cocycle, so $\theta=-\tfrac13$; then Theorem \ref{thm:c10-brst} and Corollary \ref{cor:c10-beisert}. (iii) Corollary \ref{cor:c11-theta} (iii). The converse reverses the two computations of (i).
\end{proof}

\begin{proposition}[The arity-two equation fixes each sector coefficient]
\label{prop:c11-theta-sector}
Assume the hypotheses of Theorem \ref{thm:c11-arity2} and of Definition \ref{def:c10-hyp}. Let $\mathfrak F_\perp$, $\mathfrak F_\lambda$, $\mathfrak F_V$ be families of gauge-invariant cocycles whose two-field parts have letters of the sector of $A|_{W_{\mathbb R}^\perp}$, of $\lambda$ and of $A|_V$ alone, each containing a cocycle $O$ for which $\{S^{(1)}(P_b),O\}$ is not exact. If (b) holds on $\mathfrak F_\perp$, then $\theta=-\tfrac13$; if it holds on $\mathfrak F_\lambda$, then $\theta_\lambda=-\tfrac13$; if it holds on $\mathfrak F_V$, then $\theta_V=-\tfrac13$. In particular (b) on the three families gives $\theta=\theta_\lambda=\theta_V=-\tfrac13$, the equality asserted in (a), and the cubic part of \eqref{eq:c9-arity1} admits the solution $\theta-\theta_V=\theta_\lambda-\theta_V=0$.
\end{proposition}

\begin{proof}
By Corollary \ref{cor:c11-D1local}, Proposition \ref{prop:c11-position} (i) and Lemma \ref{lem:c11-bracketpos}, the position part of $D_1$ is twice the sum of the composite pair kernels along $P_b$ minus the bracket with the position part of $\hbar F_1^{(1)}(\Delta)$. On the two-field part $\hbar F_1^{(1)}(\Delta)$ acts through its quadratic part; the bracket of the sector part of the sector of $X$ has the position part $2\partial_bX$ and the other sector parts have bracket $0$ with $X$, and the $\kappa$-directions are $s$-exact (Proposition \ref{prop:c11-arity1} (ii)). With (b), on the cohomology of the family of the sector with coefficient $\theta_\star\in\{\theta,\theta_\lambda,\theta_V\}$,
\[
\big(\text{position part of }D_1\text{ along }b\big)\cong-\tfrac23\,\text{units}\cdot\partial_b-2\theta_\star\,\text{units}\cdot\partial_b=-2\big(\tfrac13+\theta_\star\big)\,\text{units}\cdot\partial_b .
\]
The left side is $s$-exact on cocycles (Theorem \ref{thm:c11-arity2}) and $\partial_bO$ is not exact, so $\theta_\star=-\tfrac13$. The last statement is Proposition \ref{prop:c11-arity1} (iii); the existence of the quantisation is a hypothesis of Theorem \ref{thm:c11-arity2}.
\end{proof}

\begin{corollary}[The dilatation component without the cubic equation]
\label{cor:c11-noA}
Under the hypotheses of Proposition \ref{prop:c11-theta-sector}, if (b) holds on $\mathfrak F_\perp$, $\mathfrak F_\lambda$ and $\mathfrak F_V$, and (b) and (c) hold on $\mathfrak F$, then on the cohomology of $\mathfrak F$, $D_1$ induces four times the map of Theorem \ref{thm:c10-brst}, with spectrum $-8$ units times the eigenvalues of $\sum_lH_{l,l+1}$ on planar families, and $P_1\cong0$ in the translation-covariant scheme.
\end{corollary}

\begin{proof}
Proposition \ref{prop:c11-theta-sector} gives $\theta=\theta_\lambda=\theta_V=-\tfrac13$, and Theorem \ref{thm:c11-P7} (i)--(iii) applies.
\end{proof}

\begin{proposition}[The planar scalar sector]
\label{prop:c11-scalar}
On the part with the outputs $A|_{W_{\mathbb R}^\perp}$ of the planar scalar sector (Proposition \ref{prop:c10-scalarsector}), (b) and (c) hold as identities modulo terms whose sum over the cyclically adjacent pairs of every trace vanishes: per adjacent pair the composite pair kernel along $\Delta$ is $\big[-\tfrac16v.(X_aX_b)-\tfrac23X_aX_b\big]$ units and along $P_b$ it is $-\tfrac16$ units times $\partial_b(X_aX_b)$, so that on $O=\operatorname{tr}(\iota_{w_1}A\cdots\iota_{w_L}A)$ the composite term along $\Delta$ is $-\big(\tfrac13v.+\tfrac23L\big)O$ units, the composite term along $P_b$ is $-\tfrac13$ units times $\partial_bO$, and the sum of the translation-invariant parts is $-\tfrac43L$ units, $-\tfrac23$ units times the letter count. On the bilinears these are exact identities (Example \ref{ex:c11-bilinear}).
\end{proposition}

\begin{proof}
By Proposition \ref{prop:c10-scalarsector} (i), (e), the composite term of an adjacent pair is $-v.(X_aX_b)$ units, and the bracket with the renormalised composite two-point functional on a letter is $(\tfrac53v.X-\tfrac23X)$ units (Remark \ref{rem:c9-Pi}); by Theorem \ref{thm:c11-local} (v) their sum is the sum of the composite pair kernels, and the one-body term distributed over the two adjacent pairs of each letter gives the density per pair. For $L$ letters without derivatives the letter count is $2L$ and the translation-invariant part of $v.$ is $2L$: $-\tfrac13\cdot2L-\tfrac23L=-\tfrac43L$. Along $P_b$: Theorem \ref{thm:c9-P1}, $-2+\tfrac53=-\tfrac13$.
\end{proof}

